\begin{lemma}[Parameters for the partial-embedding construction]
Fix $\Delta\in\mathbb N$ and $0<d\leq1$, and put $a=d(d/2)^\Delta$.
There is $0<d'_0<d$ such that for every $0<d'\leq d'_0$ there is
$0<\varepsilon_0<\infty$ such that for every
$0<\varepsilon\leq\varepsilon_0$ there is $1\leq C_0<\infty$
such that for every $C_0\leq C<\infty$ there is $n_0\in\mathbb N$
with the following inequalities for every integer $n\geq n_0$:
\[
\begin{gathered}
 n>0,\quad \varepsilon\leq\tfrac12,\quad\varepsilon\leq d/2,
 \quad 0<a\leq d,\\
 a\leq d(d-\varepsilon)_+^k\quad(0\leq k\leq\Delta),\\
 \varepsilon(1+\varepsilon)<a,\qquad
 (\Delta+1)\varepsilon(1+\varepsilon)\leq a/2,\\
 d'\leq a/4,\qquad 2/a\leq C,\qquad 2\leq C,
 \qquad C/n\leq1.
\end{gathered}
\]
In addition, every natural number $t\geq(a/2)n$ satisfies
$1/t\leq C/n$, and every natural number $t\geq(1-\varepsilon)_+n$
is positive and satisfies $1/t\leq C/n$.
All scalars are extended nonnegative reals; the displayed bounds force
them to be finite, and the subtraction is truncated subtraction.
\end{lemma}
\begin{proof}
The assumptions make $d$ finite. Put $a=d(d/2)^\Delta$ and $d'_0=a/4$.
Since $0<d/2\leq1/2$, its $\Delta$th power is positive and at most
one, including when $\Delta=0$. Thus $0<a\leq d\leq1$ and
$0<a/4<d$. Fix $0<d'\leq a/4$ and choose
\[
 \varepsilon_0=\min\{1/2,d/2,d'/4,a/(8(\Delta+1))\}.
\]
All four entries are positive and finite, so the same holds for their
minimum. For $0<\varepsilon\leq\varepsilon_0$, choose
\[
 C_0=\max\{1,2/a,2/d',1/\varepsilon\}.
\]
Its entries are finite and it is at least one. Given finite $C\geq C_0$,
choose an integer $n_0\geq\max\{1,C\}$, using the Archimedean
property on the finite real representative of $C$. Then $n\geq n_0$
implies $n>0$, $C\leq n$, and $C/n\leq1$.
Because $\varepsilon\leq d/2$, the truncated difference is the ordinary
nonnegative difference and satisfies $d-\varepsilon\geq d/2$.
For $k\leq\Delta$ we have
$(d-\varepsilon)^k\geq(d/2)^k\geq(d/2)^\Delta$; the second
inequality uses $0<d/2\leq1$. Multiplication by $d$ gives the power
bound. Also $1+\varepsilon\leq3/2$ and
$(\Delta+1)\varepsilon\leq a/8$, so
\[
 (\Delta+1)\varepsilon(1+\varepsilon)\leq3a/16\leq a/2.
\]
Since $\Delta+1\geq1$, the same bound $3a/16$ holds for
$\varepsilon(1+\varepsilon)$, and it is strictly less than $a$
because $a>0$. Finally $d'\leq a/4$ was imposed, $C\geq2/a$ follows
from $C\geq C_0$, and $a\leq1$ gives $2/a\geq2$.
For $t\geq(a/2)n$, the lower bound is positive, so reciprocation gives
$1/t\leq2/(an)\leq C/n$. For $t\geq(1-\varepsilon)n$, we have
$t\geq n/2>0$ because $\varepsilon\leq1/2$; consequently
$1/t\leq2/n\leq C/n$.
Every computation is on finite nonnegative real representatives, where
the field operations agree with extended nonnegative real operations
because every denominator used is positive. The choices depend only
on the parameters quantified earlier, never on a graph or on $r$.
\end{proof}
